% ===================================================================
%  BAGAN No. 11 — Kota dan bangunannya. --- Kebakaran.
%
%  Traduction, faite en 2026, en indonesien d'aujourd'hui, sur l'ido
%  de texto/io. Regles de la colonne : voir 10-bagan-01.tex. Deux
%  parties, deux numerotations : les cles portent c1 et c2. Les blocs
%  t11-c1-12-1 et t11-c2-03-1 tiennent DEUX alineas sous une seule
%  cle — ordo.py les marque « ¶2 ».
%
%  IL Y A CINQ COUCHES DANS CETTE LANGUE, ET NON QUATRE. C'est la
%  DEUXIEME fois que ce compte doit etre corrige, et cette fois il
%  l'est par le violon. L'ido dit « violino » (81) ; l'indonesien dit
%  biola, et biola vient du portugais viola. En tirant ce fil on
%  s'apercoit que la colonne emploie depuis le tableau 1 vingt-trois
%  mots d'origine portugaise, cent soixante-deux fois en tout, sans
%  qu'aucun en-tete l'ait releve :
%
%    meja < mesa        jendela < janela    gereja < igreja
%    keju < queijo      sepatu < sapato     bendera < bandeira
%    mentega < manteiga garpu < garfo       kemeja < camisa
%    minggu < domingo   lemari < armário    bangku < banco
%    boneka < boneca    roda < roda         tenda < tenda
%    pesta < festa      dansa < dança       nona < dona
%    beranda < varanda  pita < fita         bola < bola
%    kereta < carreta   sekolah < escola
%
%  LE PORTUGAIS EST LA COUCHE LA PLUS ANCIENNE ET LA PLUS INVISIBLE.
%  Malacca tombe en 1511, un siecle avant les Hollandais ; les mots
%  entres alors ont eu quatre cents ans pour se fondre, et ils
%  nomment les choses les plus quotidiennes qui soient — la table, la
%  fenetre, l'eglise, la chaussure, le dimanche. Aucun locuteur ne les
%  sent etrangers, et c'est exactement pour cela qu'on ne les avait
%  pas vus : une couche se remarque tant qu'elle detonne, et celle-ci
%  ne detonne plus.
%
%  ET LA FAUTE ETAIT DEJA ECRITE DANS CETTE COLONNE. L'en-tete du
%  tableau 4 disait, a propos de beranda, « qui lui vient du portugais
%  varanda ». La preuve etait sous la main depuis six tableaux et on
%  n'en avait pas tire la consequence. Le compte des couches ne se
%  corrige pas en comptant mieux : il se corrige en relisant ce qu'on
%  a deja ecrit.
%
%  LE RESTE DU TABLEAU EST ADMINISTRATIF, ET LA C'EST LE NEERLANDAIS
%  QUI TIENT : kantor le bureau, bea cukai la douane (avec un mot
%  malais), bursa, bank, trotoar, kios, ambulans, komisaris,
%  perwira. Mais tandu (54), le brancard, est malais de fond — c'est
%  l'objet qu'on porte a la main, et le tableau 2 avait deja montre
%  que ce qu'on tient ne s'emprunte pas.
% ===================================================================

%%K t11-apar-1 apar
\VUsaut{2.0mm}
\VUcentre{13.2pt}{90}{SERI KETIGA}
\VUsaut{3.0mm}
\VUfilet{16mm}
\VUsaut{5.6mm}
\VUtitre{15.0pt}{60}{BAGAN No. 11}
\VUsaut{1.4mm}
\VUpk{11.4pt}{40}{Kota dan bangunannya. --- Kebakaran.}
\VUsaut{2.2mm}

%%K t11-tit-1 sub
\VUpk{10.2pt}{40}{Daerah pinggiran.}
\VUsaut{3.0mm}

%%K t11-c1-01-1 p
1. --- Saya tinggal di sebuah kota besar yang terletak seratus
kilometer dari samudra, dan yang toh menjadi \VUgras{pelabuhan}
\textsuperscript{(1)} kelas satu. Sungai yang membatasinya lebarnya
empat ratus meter dan membentuk kelokan yang sangat indah.

%%K t11-c1-02-1 p
2. --- Seluruh bagian kota yang berada di \VUgras{kade}
\textsuperscript{(2)} tampak sepenuhnya bercorak laut dan industri, dan
membentuk sebuah \VUgras{daerah pinggiran} \textsuperscript{(3)} yang
hanya dihuni pelaut dan buruh. Mereka bekerja di \VUgras{pabrik-pabrik}
\textsuperscript{(4)} dan di \VUgras{pabrik gas} \textsuperscript{(5)},
yang \VUgras{cerobongnya} \textsuperscript{(6)} yang tinggi dan
\VUgras{tangki gasnya} terlihat dari sangat jauh.

%%K t11-c1-03-1 p
3. --- Pada Abad Pertengahan kota itu berbenteng, dan orang masih
menemukan beberapa sisa temboknya — terutama \VUgras{gerbang}
\textsuperscript{(8)} \VUgras{kastil berbenteng} \textsuperscript{(7)}
tua, yang diapit kedua \VUgras{menaranya} \textsuperscript{(9)}, yang
kini dipakai sebagai \VUgras{penjara.}

%%K t11-c1-04-1 p
4. --- Di seluruh \VUgras{kawasan} ini, yang tampak sangat tua, hanya
satu bangunan yang agak modern : \VUgras{kantor bea cukai}
\textsuperscript{(10)}. Ia terletak di antara kade dan \VUgras{gang}
\textsuperscript{(11)} yang menuju \VUgras{balai kota}
\textsuperscript{(12)} lama. Bangunan yang terakhir ini mudah dikenali
dari \VUgras{menara loncengnya} \textsuperscript{(13)} yang raksasa,
yang menguasai kota kuno itu.

%%K t11-c1-05-1 p
5. --- Di seberang jalan \VUgras{berdiri} sebuah gedung yang dibangun
dengan gaya yang sama seperti kantor bea cukai : itulah \VUgras{bursa}
\textsuperscript{(14)}. Salah satu mukanya menghadap sebuah alun-alun
kecil tempat \VUgras{kantor bupati} \textsuperscript{(15)} berada.
Sesungguhnya kota kami, tanpa menjadi ibu kota negara, tetap menjadi
ibu kota daerah yang penting.

%%K t11-tit-2 sub
\VUsaut{5.0mm}
\VUpk{10.2pt}{40}{Bagian kota yang paling tinggi.}
\VUsaut{3.4mm}

%%K t11-c1-06-1 p
6. --- Garnisun, yang cukup besar jumlahnya, ditempatkan di
\VUgras{barak-barak} \textsuperscript{(16)} yang luas dan sangat sehat ;
barak-barak itu terbentang sampai ke pagar \VUgras{taman kota}
\textsuperscript{(17)}. \VUgras{Jalan besar} \textsuperscript{(19)} yang
menuju taman itu lewat di belakang \VUgras{bank tabungan}
\textsuperscript{(18)} dan berakhir di alun-alun \VUgras{katedral}
\textsuperscript{(20)}.

%%K t11-c1-07-1 p
7. --- Semua bangunan itu terbentang seperti amfiteater di atas sebuah
bukit kecil, tempat \VUgras{sekolah menengah} \textsuperscript{(21)}
kami yang luas juga berada.

%%K t11-c1-07-2 p
Bila orang turun lewat jalan yang agak menurun, yang menyambung dengan
Jalan Besar, ia sampai di \VUgras{pengadilan} \textsuperscript{(22)},
yang di depannya terbentang sebuah \VUgras{taman umum}
\textsuperscript{(26)} yang menyenangkan. \VUgras{Taman kecil} itu
tidak besar, tetapi ia membuat di tengah kota sebuah oase kehijauan dan
kesejukan. Karena itu ia sangat sering dikunjungi penduduk kota, yang
senang datang duduk di bawah tenda \VUgras{kedai kopi}
\textsuperscript{(25)} untuk mendengarkan para pemusik tentara yang
bermain pada hari Kamis dan hari Minggu di \VUgras{panggung musik}
\textsuperscript{(27)} yang diletakkan di antara petak rumput dan
rumpun bunga.

%%K t11-c1-08-1 p
8. --- Di salah satu petak rumput itu berdiri sebuah \VUgras{patung}
\textsuperscript{(28)} perunggu, tugu yang didirikan untuk mengenang
salah seorang warga kota kami, yang telah memberikan jasa besar kepada
kota kelahirannya.

%%K t11-tit-3 sub
\VUsaut{4.0mm}
\VUpk{10.2pt}{40}{Sarana angkutan.}
\VUsaut{3.4mm}

%%K t11-c1-09-1 p
9. --- Sampai sekarang kota kami belum diterangi dengan cahaya
listrik ; ia hanya mempunyai \VUgras{lampu-lampu gas}
\textsuperscript{(29)}. Tetapi \VUgras{surat kabar di sini} berkampanye
supaya sistem penerangan itu diperbaiki, \VUgras{seperti} orang telah
menyempurnakan \VUgras{sistem penggerak trem.} \VUgras{Kereta-kereta}
tua yang \VUgras{ditarik} \VUgras{kuda} praktis diganti dengan
\VUgras{trem listrik} \textsuperscript{(31)}, yang mengangkut penumpang
dengan cepat dari \VUgras{ujung kota yang satu ke ujung yang lain}
dengan harga yang sangat \VUgras{murah.}

%%K t11-c1-10-1 p
10. --- Di dekat gedung pengadilan ada \VUgras{bank}
\textsuperscript{(32)} tempat orang mendiskontokan surat berharga
dagang. \VUgras{Para penagih bank} \textsuperscript{(33)} pergi menagih
utang ke rumah-rumah.

%%K t11-tit-4 sub
\VUsaut{4.0mm}
\VUpk{10.2pt}{40}{Seni dan pendidikan.}
\VUsaut{3.4mm}

%%K t11-c1-11-1 p
11. --- Gedung indah yang berdiri di dekatnya adalah \VUgras{museum.}
Ia sekaligus berisi \VUgras{perpustakaan} \textsuperscript{(34)} kota,
\VUgras{koleksi ilmu alam} \textsuperscript{(35)} yang bagus (Museum
Alam) dan sebuah \VUgras{galeri lukisan} \textsuperscript{(36)} yang
anggun, yang merupakan \VUgras{ruang pameran} tetap yang megah.

%%K t11-c1-11-2 p
Ia dibuka untuk umum dua kali seminggu, tetapi orang asing boleh
mengunjunginya pada hari apa saja dengan memberi sedikit uang kepada
\VUgras{penjaga} \textsuperscript{(37)}. \VUgras{Para pelukis}
\textsuperscript{(38)} datang bekerja di sana. Orang melihat mereka
\VUgras{di depan} kuda-kuda mereka, dengan \VUgras{palet} di tangan,
menyalin lukisan-lukisan yang termasyhur.

%%K t11-c1-12-1 p
12. --- Di ketinggian, di belakang museum, mulai tampak \VUgras{kubah}
\textsuperscript{(40)} \VUgras{rumah sakit} \textsuperscript{(39)} dan
bangunan-bangunan \VUgras{sekolah guru} \textsuperscript{(41)} tempat
para guru sekolah kota kami dilatih.

Di alun-alun Teater ada \VUgras{kantor pos} \textsuperscript{(43)} yang
dipasang di dalam sebuah \VUgras{rumah pribadi}
\textsuperscript{(42)}.

%%K t11-tit-5 sub
\VUsaut{5.0mm}
\VUpk{11.4pt}{40}{Kebakaran.}
\VUsaut{3.0mm}

%%K t11-tit-6 sub
\VUpk{10.2pt}{40}{Teriakan kebakaran.}
\VUsaut{3.4mm}

%%K t11-c2-01-1 p
1. --- Kabar yang mengerikan tersebar ke seluruh kota : kebakaran baru
saja terjadi di gedung teater! Pada saat saya sampai di
\VUgras{alun-alun} \textsuperscript{(96)}, kerumunan sudah berkumpul di
\VUgras{trotoar} \textsuperscript{(46)} dan di \VUgras{jalan raya}
\textsuperscript{(47)}. \VUgras{Para pegawai pos} \textsuperscript{(44)}
dan \VUgras{para pengantar surat} \textsuperscript{(45)} keluar dari
kantor pos. Seorang \VUgras{sopir trem} \textsuperscript{(48)} harus
menghentikan keretanya untuk memberi jalan kepada \VUgras{mobil
ambulans} \textsuperscript{(50)} kota, yang datang membawa seorang
\VUgras{ahli bedah} \textsuperscript{(52)} dan seorang \VUgras{perawat}
\textsuperscript{(51)} untuk merawat seorang \VUgras{yang luka}
\textsuperscript{(53)}, yang dibawa dengan \VUgras{tandu}
\textsuperscript{(54)} ke dekat \VUgras{kios koran}
\textsuperscript{(55)}.

%%K t11-c2-02-1 p
2. --- Sebuah kompi infanteri, yang sedang kembali dari latihan,
berhenti untuk menjamin ketertiban bersama \VUgras{polisi kota
berkuda} \textsuperscript{(61)}. \VUgras{Para penunggang,} yang kudanya
berdiri di kaki belakang, lebih berhasil daripada \VUgras{para
prajurit} \textsuperscript{(57)} atau \VUgras{polisi militer}
\textsuperscript{(63)} dalam mendorong mundur \VUgras{orang-orang yang
ingin tahu} \textsuperscript{(56)}. Lagi pula polisi militer sangat
sibuk. Salah seorang dari mereka menunjukkan kepada \VUgras{polisi}
\textsuperscript{(64)} ke mana harus dibawa seorang \VUgras{korban
kecelakaan} \textsuperscript{(65)} yang lain, yaitu seorang pemadam
kebakaran yang berani yang jatuh dari tangga.

%%K t11-c2-03-1 p
3. --- \VUgras{Perwira} \textsuperscript{(66)} menentukan tempat
\VUgras{pompa uap} \textsuperscript{(67)} yang sedang datang harus
dipasang. Sambil menunggu mesin yang perkasa itu, ia dengan sigap
memasang sebuah \VUgras{pompa tangan} \textsuperscript{(68)} yang
\VUgras{selangnya} \textsuperscript{(69)} yang panjang menyemburkan air
seperti arus deras ke atas api.

Enam pemadam kebakaran menggerakkannya, sementara kawan mereka, yang
memegang \VUgras{mulut selang} \textsuperscript{(70)}, mengarahkan
\VUgras{pancaran} \textsuperscript{(71)} ke pusat kebakaran.

%%K t11-c2-04-1 p
4. --- Di sisi lain gedung teater, orang menyandarkan \VUgras{tangga}
\textsuperscript{(72)} yang besar. Tangga itu memungkinkan para pemadam
kebakaran mendekati lidah api yang menyembur di antara pusaran
\VUgras{asap} \textsuperscript{(75)} hitam.

%%K t11-tit-7 sub
\VUsaut{5.0mm}
\VUpk{10.2pt}{40}{Perbuatan berani.}
\VUsaut{3.4mm}

%%K t11-c2-05-1 p
5. --- \VUgras{Kebakaran} \textsuperscript{(74)} itu lahir di ruang
depan \VUgras{gedung teater} \textsuperscript{(73)} ketika orang sedang
melatih sebuah \VUgras{opera} yang \VUgras{pertunjukan} pertamanya akan
berlangsung besok. Untunglah hanya sedikit \VUgras{penonton}
\textsuperscript{(76)} yang berada di ruang pertunjukan. Orang-orang
malang itu berlindung di \VUgras{balkon} \textsuperscript{(77)}, tempat
\VUgras{para pemadam kebakaran} \textsuperscript{(86)} yang berani
datang menolong mereka dengan tangga dan tali besar.

%%K t11-c2-06-1 p
6. --- Di bawah \VUgras{serambi bertiang} \textsuperscript{(78)},
tempat \VUgras{poster} \textsuperscript{(79)} memberitakan pertunjukan
itu, orang melihat \VUgras{para pemain musik} \textsuperscript{(81)}
\VUgras{orkes} \textsuperscript{(80)} melarikan diri sambil membawa
alat-alat mereka : \VUgras{biola} \textsuperscript{(81)},
\VUgras{suling} \textsuperscript{(82)}, \VUgras{selo}
\textsuperscript{(83)}, \VUgras{trombon} \textsuperscript{(84)},
\VUgras{tambur} \textsuperscript{(85)}, dan sebagainya. Orang berusaha
menyelamatkan sebagian \VUgras{perabot} \textsuperscript{(87)}.

%%K t11-c2-07-1 p
7. --- \VUgras{Para aktor} \textsuperscript{(89)} dan \VUgras{para
aktris} \textsuperscript{(88)}, \VUgras{para penyanyi laki-laki} dan
\VUgras{perempuan} terpaksa melarikan diri dengan \VUgras{kostum}
\textsuperscript{(90)} panggung mereka. Penyanyi tenor menjelaskan
kepada seorang \VUgras{wartawan} \textsuperscript{(92)} bagaimana hal
itu terjadi. \VUgras{Pelapor} itu segera berlari datang sejak saat
pertama bersama para pejabat : \VUgras{bupati} \textsuperscript{(91)},
\VUgras{wali kota} \textsuperscript{(93)}, \VUgras{komisaris polisi}
\textsuperscript{(94)} dan \VUgras{pejabat pengadilan}
\textsuperscript{(95)}, yang akan mengadakan penyelidikan untuk
memutuskan siapa yang bertanggung jawab.

\VUsaut{6.0mm}
\VUfilet{22mm}
